Pragmatic sentiment in Vietnamese social-media text involves more than surface polarity: an utterance may convey implicit sentiment, sarcasm, irony, idiom/figurative language, code-switching, or mocking. We introduce the ViPragSent dataset and an XLM-R-large multi-task model that predicts these six binary pragmatic phenomena while using three-way polarity and seven-way emotion as separate auxiliary tasks. The proposed model also uses rationale supervision only during training; inference uses classification heads. On the in-domain test cohort, XLM-R-large multi-task (ours) obtains $93.67 \pm 0.19$ macro-pragmatic F1 and the highest observed mean on each pragmatic head among the evaluated systems with complete three-run results. We further examine external affective transfer, auxiliary-objective trade-offs, positive sarcasm supervision budgets, and probability calibration. The analyses separate in-domain discrimination from transfer, operating cost, and calibration behavior. Rationale supervision is training-only and is not part of the inference interface.
